\sectionlink{単一場インフレーションのさまざまな極限}

\subsectionlink{スローロール・インフレーションと高エネルギー補正}
一般的なラグランジアン (\ref{eq:actiontad}) の最も単純な例は，背景のHubbleパラメータ $H(t)$ がわかれば決まる最初の三つの項だけを残し，それ以外の高次の演算子をすべてゼロ，すなわち $M_2 = M_3 = \bar M_1 =\bar M_2 \ldots =0$ とすることで得られる．$\phi$ を用いた記述では，これは高次の項を含まない標準的なスローロール・インフレーションに対応する．
この場合，前節で論じたように，重力との混合を無視してGoldstoneモードのラグランジアン (\ref{Spi}) に着目することで，スケール $H$ における予言を行える．計算したいのは，地平線通過の直後における保存量 $\zeta$ である．これは線形の次数では，ゲージを $\pi = 0$ （本論文の用語ではユニタリーゲージ）に選び，計量の空間成分を
\be
\label{eq:zetagauge}
g_{ij} = a^2(t) \left[(1+ 2 \zeta(t,\vec x)) \delta_{ij} + \gamma_{ij}\right]
\ee
とすることで定義される．ここで $\gamma$ は横波かつトレースレスであり，二つのテンソル自由度を記述する．
$\pi$ と $\zeta$ の関係は非常に単純である．重力との混合を無視しているので，$\pi$ を用いた記述では計量に摂動はない．$\pi = 0$ とするには時間微分同相変換 $t \to t - \pi(t,\vec x)$ を行えばよく，これによって (\ref{eq:zetagauge}) の形の空間計量が得られ，
\be
\label{eq:pizeta}
\zeta (t,\vec x) = - H \pi(t, \vec x) \;.
\ee
となる．

各モード $k$ について，関心があるのは地平線通過 $\omega(k) = k/a \sim H$ の前後の力学だけである．この期間には，スローロール補正を除いて背景をde Sitter時空で近似できる．したがって，正準規格化されたスカラー $\pi_c$ の2点関数は，de Sitter時空での結果
\be
\langle \pi_c(\vec k_1) \pi_c(\vec k_2)\rangle = (2 \pi)^3 \delta(\vec k_1 + \vec k_2) \frac{H^2_*}{2 k_1^3} \;,
\ee
で与えられる．ここでも以下でも，$*$ は地平線通過時の量の値を表す．これから，$\zeta$ の2点関数は
\be
\label{eq:spectrumzeta}
\langle \zeta(\vec k_1) \zeta(\vec k_2)\rangle = (2 \pi)^3 \delta(\vec k_1 + \vec k_2) \frac{H^4_*}{4 M_{\rm Pl}^2 |\dot H_*|}\frac{1}{k_1^3} = (2 \pi)^3 \delta(\vec k_1 + \vec k_2) \frac{H_*^2}{4 \epsilon_* M_{\rm Pl}^2} \frac{1}{k_1^3} \;.
\ee
で与えられることがわかる．変数 $\zeta$ は地平線の外側で一定なので，この式はスローロール補正を除いてすべての $k$ に対して厳密である．特に，これを用いればスローロールの最低次でスペクトルの傾きを計算でき，
\be
n_s-1= \frac{d}{d \log k} \log \frac{H_*^4}{|\dot H_*|} = \frac1{H_*} \frac{d}{d t_*} \log \frac{H_*^4}{|\dot H_*|} = 4 \frac{\dot H_*}{H_*^2} - \frac{\ddot H_*}{H_* \dot H_*} \;.
\ee
となる．

ただし，すべての観測量が $\pi$ のラグランジアン (\ref{Spi}) から計算できるわけではないことに注意しよう．最低次の結果が重力との混合に由来する場合，あるいはスローロール展開の高次で生じる場合には，この方法では計算できない．例えば，式~(\ref{Spi}) の最初の二つの項は $\pi$ の自己相互作用を含まないため，3点関数 $\langle \zeta(\vec k_1) \zeta(\vec k_2) \zeta(\vec k_3) \rangle $ はゼロになってしまう．したがって，式~(\ref{Smixed}) における重力との混合を考慮して，次の次数の補正を調べる必要がある．

高次の項をすべてゼロにするという選択が，厳密には成り立ち得ないことは明らかである．ツリーレベルでゼロにしたとしても，少なくとも輻射補正によってそれらは生成される．この理論は繰り込み不可能であり，摂動展開を十分高次まで進めると，あらゆる相互作用が発散する係数を伴って生成される．追加の項は重力子のループによって生成されるので，非常に小さいこともあり得る．例えば，ユニタリーゲージの相互作用 $M_{\rm Pl}^2 \dot H g^{00}$ から出発すると，$(g^{00}+1)^2$ の形の項が，対数的に発散する係数 $M_2^4 \sim \dot H^2 \log
\Lambda$ を伴って生成されることは容易に確かめられる．これは，$M^4_2 \gtrsim \dot H^2$ と仮定すべきことを意味する（\footnote{重力に結合した質量ゼロのスカラー場の理論における対数発散の具体的な計算は，かなり以前に文献 \cite{tHooft:1974toh} で行われている．}）．ただし，この下限は非常に小さい．例えば，$\pi$ の分散関係は，時間方向の運動項への追加の寄与によって変化する．以下で詳しく論じるように，これは $\pi$ の励起の速度が光速からわずかにずれ，その相対的なずれが $1-c_s \sim M_2^4/(|\dot H| M_{\rm Pl}^2) \sim |\dot H|/M_{\rm Pl}^2$ となることを意味する．スカラーのスペクトルの規格化，式~(\ref{eq:spectrumzeta}) を用いると，光速からのずれは $\gtrsim \epsilon^2 \cdot 10^{-10}$ であることがわかる．あまり興味深い下限ではない．

Planckスケールよりも低いスケールで新たな物理が加わるなら，追加の演算子の大きさはずっと大きくなる．一般に，本論文の手法は，新しい物理がもたらし得るあらゆる効果を正しくパラメータ化する．有効場理論の手法では通常そうであるように，模型の紫外完備化の詳細は高次元演算子に符号化される．これは，標準模型を超える物理で起こることと非常によく似ている．低エネルギーでは，新しい物理がもたらし得る効果は，対称性と両立する一連の高次元演算子に符号化される \cite{Barbieri:2004qk}．標準模型を実験で詳しく調べることによって，これらの高次元演算子の大きさに厳しい制限を課すことができる．本論文の場合にも同じことができるが，残念ながら，考えられる観測の種類ははるかに限られる．
高次元演算子に対して課し得る実験的制限の一例は，重力波のスペクトルの傾きに関する整合性関係である．よく知られているように，Einstein--Hilbert作用から得られる重力波のスペクトルは，
\be
\label{eq:GWspectrum}
\langle \gamma^s(\vec k_1) \gamma^{s'}(\vec k_2)\rangle = (2 \pi)^3
\delta(\vec k_1 + \vec k_2) \frac{H_*^2}{M_{\rm Pl}^2}\frac{1}{k_1^3}
\delta_{s s'}\,
\ee
で与えられる．ここで $\gamma^s$ は，重力波の二つの偏極を表す．この寄与とスカラーの寄与 (\ref{eq:spectrumzeta}) の比は $\epsilon_*$ で与えられる．したがって，重力波のスペクトルの傾き $n_g=-2 \epsilon_*$ は，テンソルモードとスカラーモードの比がわかれば決まる．

この予言は，$M_2 =0$，{\em すなわち} $c_s=1$ を仮定する場合に有効である．実際，後で見るように，スカラーのスペクトルは $c_s^{-1}$ に比例する一方，重力波に関する予言は $M_2$ によって変わらない．整合性関係が実験的に検証されれば，たとえ誤差が大きくても，$c_s$ が $1$ から大きくずれ得ないことがわかり，これは
\be
M_2^4 \lesssim M_{\rm Pl}^2 |\dot H| \;.
\ee
を意味する．

高次元演算子はスカラーの揺らぎだけでなくテンソルモードにも影響することに注意しよう．もっとも，これらの補正を検証するのは，おそらくはるかに難しい．例えば，スカラーの揺らぎとの関連を後で論じるユニタリーゲージの演算子 $-\bar M_3(t)/2  \cdot \delta K^\mu {}_\nu \delta
K^\nu {}_\mu $ は，$\dot g_{ij}^2$ の形の項を含む．これは重力波の分散関係を変える．実際，この演算子が存在するときのテンソルモード $\gamma_{ij}$ の作用は容易に求められ，
\be
S_{\gamma}=\frac{M_{\rm Pl}^2}{8} \int d^4x \sqrt{-g} \left[\left(1-\frac{\bar M_3^2}{M_{\rm Pl}^2}\right)\dot\gamma_{ij} \dot\gamma_{ij}- \frac1{a^2}
\partial_l\gamma_{ij} \partial_l\gamma_{ij} \right] \;.
\ee
を得る．したがって，重力波のスペクトル (\ref{eq:GWspectrum}) は $\bar M_3^2/M_{\rm Pl}^2$ 程度の補正を受ける．$\bar M_3^2$ をPlanckスケールまで大きくしない限り，この補正は小さい．$(g^{00}+1)^{n}$ の形の演算子は，計量の横波・トレースレス成分に影響しないため，テンソルモードに影響しないことは容易にわかる．

さまざまな演算子に対する実験的制限の他の例については，以下の節で論じる．

\subsectionlink{小さな音速と大きな非Gaussian性}

Goldstoneモードの作用 (\ref{Spi}) から，空間方向の運動項 $(\partial_i \pi)^2$ は背景の時間発展によって完全に決まり，$M_{\rm Pl}^2 \dot H (\partial_i\pi)^2$ となることがわかる．特に，$\dot H <0$ の場合にのみ，この項は「健全な」負の符号を持つ．これは，ヌル・エネルギー条件の破れと系の不安定性の存在との間の，よく研究されてきた関係の一例である．FRW宇宙では，この条件の破れは $\dot H<0$ と同値である \cite{Hsu:2004vr,Dubovsky:2005xd}．ただし，演算子 $(\partial_i \pi)^2$ の符号が不適切であるというだけでは，系に病的な問題があるとは結論できないことに注意しよう．以下で詳しく論じるように，特定の領域では $\delta K^\mu {}_\mu {}^2$ のような高次の項が重要になり得る．文献 \cite{Creminelli:2006xe} では，病的な問題を生じることなく $\dot H >0$ を実現できる例が研究されている．

一方，時間方向の運動項 $\dot{\pi}^2$ の係数は，二次の演算子 $(g^{00}+1)^2$ からも寄与を受けるため，背景の時間発展だけでは完全には決まらない．式~(\ref{Spi}) では，
\begin{equation}
\left(-M_{\rm Pl}^2 \dot{H} + 2 M_2^4 \right) \dot\pi^2 \;.
\end{equation}
となっている．不安定性を避けるには，$-M_{\rm Pl}^2 \dot{H} + 2 M_2^4 >0$ でなければならない．
時間方向と空間方向の運動項の係数が異なるので，$\pi$ の波の「音速」は $c_s \neq 1$ となる．背景はLorentz不変性を自発的に破っており，$c_s=1$ はいかなる対称性によっても保護されないため，これは予想されることである．前節で論じたように，$c_s=1$ からのずれは，少なくとも重力子のループによって誘起される\footnote{重力との結合と，ユニタリーゲージのラグランジアンにおける演算子の時間依存性を無視すると（そのため $\pi \to \pi + {\rm const}$ が対称性となる），$c_s=1$ は対称性 $\partial_\mu\pi \to
\partial_\mu\pi + v_\mu$ によって保護できる．ここで $v_\mu$ は定数ベクトルである．この対称変換の下で，$\pi$ のLorentz不変な運動項の変化は全微分となる一方，式~(\ref{Spi}) における $M_2^4$ に比例する演算子は明らかに不変ではないため，$c_s=1$ となる．$\pi$ により多くの微分が作用する相互作用を記述することは許されるので，この理論が自由理論ではないことに注意しよう．この対称性は，DGP模型のブレーンの曲がりモードの研究に現れる \cite{Adams:2006sv}．}．
音速は
\be
c_s^{-2} = 1-\frac{2 M_2^4}{M_{\rm Pl}^2 \dot H} \;.
\ee
で与えられる．これは，超光速伝播を避けるためには $M_2^4 >0$ でなければならないことを意味する（$\dot H <0$ を仮定する）．超光速伝播は，この理論にLorentz不変な紫外完備化が存在しないことを意味する \cite{Adams:2006sv}．以下では $c_s \leq 1$ の場合に着目する．$c_s > 1$ の模型についての現象論的な議論は，文献 \cite{Mukhanov:2005bu} を参照されたい．

上の $c_s^2$ の式を用いると，Goldstoneモードの作用は三次までで

\begin{eqnarray}\label{Spi_cs}
&&S_{\rm \pi} = \int \!  d^4 x  \:  \sqrt{- g} \left[
-\frac{M^2_{\rm Pl} \dot{H}}{c^2_s} \left(\dot\pi^2-c^2_s
\frac{(\partial_i\pi)^2}{a^2}\right)
\right.\nonumber \\ &&\qquad\left.+M_{\rm Pl}^2 \dot H \left(1-\frac{1}{c^2_s}\right)
\left(\dot{\pi}^3-\dot\pi\frac{(\partial_i\pi)^2}{a^2}
\right)- \frac43 M_3^4 \dot\pi^3... \right].
\end{eqnarray}

と書ける．

\ref{sec:Goldstone}節での議論から，エネルギー $E \gg E_{\rm mix}
\simeq M_2^2/M_{\rm Pl}$ では重力との混合を無視できることがわかっている．これは，$H$ 程度のエネルギーで行われる宇宙論的観測量の予言が，$H \gg M_2^2/M_{\rm Pl}$，あるいは同値な条件 $\epsilon/c_s^2 \ll 1$ が成り立つ場合には，Goldstoneモードの作用 (\ref{Spi_cs}) によって最低次で捉えられることを意味する．そうでない場合には，Goldstoneモードの作用が主要な効果を含んでいるという保証はない．

2点関数の計算は，$c_s=1$ の場合とほぼ同様に進められる．その際，再スケールした運動量 $\bar k=c_s k$ を用い，時間方向の運動項の前に付く追加の因子 $c_s^{-2}$ を考慮すればよい．その結果，
\be
\label{eq:spectrumzetacs}
\langle \zeta(\vec k_1) \zeta(\vec k_2)\rangle = (2 \pi)^3 \delta(\vec
k_1 + \vec k_2)  \frac1{c_{s*}} \cdot \frac{H^4_*}{4 M_{\rm Pl}^2 |\dot H_*|}\frac{1}{k_1^3} = (2
\pi)^3 \delta(\vec k_1 + \vec k_2) \frac1{c_{s*}} \cdot \frac{H_*^2}{4 \epsilon_* M_{\rm Pl}^2} \frac{1}{k_1^3} \;.
\ee
を得る．音速の時間変化は，スペクトルの傾きに追加の寄与をもたらし，
\begin{equation}
n_s=\frac{d}{d \log k}\log \frac{H^4_{*}}{|\dot{H}_*|
c_{s*}}=\frac{1}{H_*}\frac{d}{dt_*}\log \frac{H^4_{*}}{|\dot{H}_*|
c_{s*}} =4\frac{\dot{H_*}}{H^2_*}-\frac{\ddot{H}_*}{\dot{H}_*
H_* }-\frac{\dot{c}_{s*}}{c_{s*}H_{*}} \;.
\end{equation}
となる．この結果は，文献 \cite{Garriga:1999vw} で得られた結果と一致する．

作用 (\ref{Spi_cs}) から，音速を低下させるのと同じ演算子が，$\dot{\pi}(\nabla\pi)^2$ および $\dot\pi^3$ の形のGoldstoneモードの三次結合を誘起することが明確にわかる．時間微分同相変換の非線形実現は，音速の低下と3点相関関数の増大，{\em すなわち} 非Gaussian性の増大との間に関係を課す．この関係は，文献 \cite{Chen:2006nt} における3点関数の具体的な計算で強調された．

非Gaussian性の大きさを見積もるには，凍結が起こる $\omega \sim H$ の前後で，非線形補正と二次の項を比較しなければならない．
$c_s\ll 1$ の極限では，演算子 $\dot\pi(\nabla{\pi})^2$ が主要な寄与を与える．二次の作用から，モードは $k
\sim H/c_s$ で凍結することがわかり，時間微分に比べて空間微分が増強されるためである．したがって，非Gaussian性の大きさは比
\begin{equation}
\frac{{\cal L}_{\dot\pi(\nabla \pi)^2}}{{\cal L}_{2}}\sim \frac{H
\pi \left(\frac{H}{c_s} \pi\right)^2}{H^2 \pi^2} \sim \frac{H}{c^2_s}
\pi \sim \frac{1}{c^2_s} \zeta \, ,
\end{equation}
で与えられる．最後の段階では，$\pi$ と $\zeta$ の間の線形関係，式~(\ref{eq:pizeta}) を用いた．$\zeta \sim
10^{-5}$ とすると，非線形補正の大きさを見積もることができる\footnote{非線形補正の大きさは，$\zeta$ の具体的な値に依存する．$\zeta$ の典型的な値が小さくても，例えば原始ブラックホールの生成を研究するために，非常に大きな（したがって非常に起こりにくい）揺らぎに関心を持つことはあり得る．$\zeta$ の値が十分に大きく，$\zeta \gtrsim c_S^2$ となると，非線形補正は1程度となり，摂動展開は破綻する．したがって，$\zeta$ の非常に大きな値に依存する予言は，有効場理論の適用範囲外となる可能性がある．}．
通常，非Gaussian性の大きさはパラメータ $f_{\rm
NL}$ で与えられ，パラメータ依存性は ${\cal L}_{\dot\pi(\nabla
\pi)^2}/{\cal L}_2\sim f_{\rm NL} \zeta$ という形になる．したがって，主要な寄与から
\begin{equation}
\label{eq:nabla2dot}
f^{\rm equil.}_{\rm NL, \;\dot\pi(\nabla\pi)^2}\sim \frac{1}{c^2_s} \;.
\end{equation}
を得る．上付き添字の ``equil.'' は3点関数の運動量依存性を表し，これらの模型では，いわゆる正三角形型になる \cite{Babich:2004gb}．Goldstoneモードを用いた記述では，これは物理的に明らかである．関与する $\pi$ の相互作用は微分を含むため地平線の外側で急速に消失し，同程度の波長を持つモードの間にしか相関がないからである．

Goldstoneモードのラグランジアン (\ref{Spi_cs}) には，3点関数に寄与する，もう一つの独立な演算子 $-\frac43 M_3^4 \dot\pi^3$ が存在する．これは，ユニタリーゲージの演算子 $(g^{00}+1)^3$ に由来する．したがって，$\dot\pi^3$ の形の寄与は二つあり，
\begin{equation}\label{fnldotpicube}
f^{\rm equil.}_{\rm NL,\ \dot\pi^3}\sim 1-\frac{4}{3}
\frac{M_3^4}{M_{\rm Pl}^2 |\dot H| c_s^{-2}} \;.
\end{equation}
を与える．演算子 $-\frac43 M_3^4 \dot\pi^3$ の大きさは，時間微分同相変換の非線形実現によって制限されず，自由なパラメータである．DBIインフレーション \cite{Alishahiha:2004eh} では，$M_3^4 \sim M_{\rm Pl}^2 |\dot H| c_s^{-4}$ となるので，非Gaussian性への寄与は式~(\ref{eq:nabla2dot}) の寄与と同程度である．$M_3^4$ は，ユニタリーゲージの演算子 $M_2^4 (g^{00}+1)^2$ と $M_3^4 (g^{00}+1)^3$ がともに同じエネルギースケールで強結合になると仮定した場合にも，およそ同じ大きさになる．

非Gaussian性への実験的制限を，ユニタリーゲージの演算子 $(g^{00}+1)^2$ の大きさ，したがって音速に対する制約として捉えることは興味深い．具体的な計算 \cite{Chen:2006nt} によって，演算子 $\dot\pi
(\nabla\pi)^2$ が，実験的に制約されているパラメータ $f_{\rm
  NL}^{\rm equil.}$ に与える寄与が得られる．$c_s^{-1}$ について支配的な項をとると\footnote{これは，文献 \cite{Chen:2006nt} の記法で $P_{,XXX}=0$ と置くことで得られる．}
\be
f_{\rm NL}^{\rm equil.} = \frac{85}{324} \cdot \frac1{c_s^2} \;.
\ee
となる．実験的に許される範囲 \cite{Creminelli:2006rz}
\be
\label{eq:fnleqexp}
-256 < f_{\rm NL}^{\rm equil.} < 332 \quad {\rm at} \; 95\% \;{\rm C.L.}
\ee
は，制約
\be
c_s > 0.028 \quad {\rm at} \; 95\% \;{\rm C.L.}
\ee
に置き換えられる．ただし，演算子 $\dot\pi (\nabla\pi)^2$ と $\dot\pi^3$ は，原理的には3点関数に異なる運動量依存性を与えるものの，現時点ではこの差を実験的に識別できないため，制約 (\ref{eq:fnleqexp}) は二つの演算子を合わせた効果に対するものであることに注意しよう．二つの演算子の間で打ち消し合いが生じない限り，音速への制約は成り立つ．例えばDBIインフレーションでは，先ほど論じたように，演算子 $M_3^4 (g^{00}+1)^3$ の効果は無視できない．しかし，打ち消し合いは生じず，音速への制約は
\be
{\rm DBI:} \quad c_s > 0.031 \quad {\rm at} \; 95\% \;{\rm C.L.}
\ee
へとわずかに変化するだけである．

ここまではGoldstoneモードのラグランジアンに着目してきたが，この一般的な手法は，重力との混合をすべて考慮したい場合にも有用であることを強調しておきたい．例えば，ユニタリーゲージのラグランジアン (\ref{eq:actiontad}) に戻れば，3点関数の計算にいくつの係数が関与するかを容易に見て取れる．スローロールと微分の最低次では，すでに論じたとおり，係数は $M_2$ と $M_3$ の二つである．スローロールの一次では，係数の緩やかな変化を記述する新たなパラメータが四つ現れる．すなわち，通常のスローロールパラメータ $\epsilon$ と $\eta$，そして演算子 $(g^{00}+1)^2$ と $(g^{00}+1)^3$ の係数に対する二つの追加のパラメータである．実際，これは具体的な計算 \cite{Chen:2006nt} で得られる結果にほかならない\footnote{具体的な計算から，係数の一つは3点関数に独立な運動量依存性を生じないため，他のパラメータから切り分けられないことがわかる．}．

これらの議論はすべて，4点関数（およびより高次の相関関数）へ容易に拡張できる．Goldstoneモードのラグランジアンでは，4点関数に寄与する演算子が三つある（ここでもスローロールと微分の最低次で考える）．それらは，$(g^{00}+1)^2$ （音速 $c_s$ によって決まる），$(g^{00}+1)^3$，$(g^{00}+1)^4$ である．音速によって決まる演算子の効果を見積もろう．3点関数の場合と同様に，その効果は演算子 $(\nabla\pi)^4$ によって支配され，それが誘起する非Gaussian性の大きさは
\be
\frac{{\cal L}_{(\nabla \pi)^4}}{{\cal L}_{2}}\sim \frac{\left(\frac{H}{c_s} \pi\right)^4}{H^2 \pi^2} \sim \frac{H^2}{c^4_s}
\pi^2 \sim \frac{1}{c^4_s} \zeta^2 \;.
\ee
と見積もれることが容易にわかる．これは，文献 \cite{Chen:2006dfn} で行われた具体的な計算と一致する．

\subsubsectionlink{カットオフと自然性}
これまで論じたように，$c_s <1$ の場合，Goldstoneモードの作用は繰り込み不可能な相互作用を含む．したがって，Goldstoneモード間の自己相互作用は，あるエネルギースケールで強結合となり，これが本理論のカットオフを定める．このカットオフは，ツリーレベルの部分波ユニタリティを調べること，{\em すなわち} $\pi$ のツリーレベルでの散乱がユニタリである最大のエネルギーを求めることで見積もれる．計算は容易であり，唯一の複雑さは非相対論的な分散関係に由来する．カットオフスケール $\Lambda$ は，
\be
\label{eq:cutoff}
\Lambda^4 \simeq 16 \pi^2 M_2^4 \frac{c_s^7}{(1-c_s^2)^2} \simeq 16
\pi^2 M_{\rm Pl}^2 |\dot H|
\frac{c_s^5}{1-c_s^2} \;.
\ee
となる．同じ結果は，$\pi \pi$ の散乱振幅へのループ補正が重要となるエネルギースケールを調べることでも得られる．予想どおり，$c_s$ が小さくなるほど理論はより強く結合し，カットオフスケールは低くなる．一方，$c_s \to 1$ ではカットオフは次第に高くなる．この極限では繰り込み不可能な相互作用がなく，カットオフをPlanckスケールまで引き上げられるので，これは理にかなっている．このカットオフスケールは，ユニタリーゲージの演算子 $(g^{00}+1)^2$ だけを調べて得られたものである．他の独立な演算子は，その大きさによっては，さらに低いエネルギーのカットオフを与える可能性がある．スケール $\Lambda$ は，本理論が弱結合であり，意味を持つ最大のエネルギーを示していることに注意しよう．このスケールよりも下で，新しい物理が登場しなければならない．ただし，新しい物理は $\Lambda$ よりもはるかに低いスケールで現れる可能性もある．

このラグランジアンを用いて宇宙論的な相関関数を予言したいなら，インフレーション中のHubbleパラメータ $H$ 程度のスケールでこれを用いる必要がある．したがって，このエネルギースケールはカットオフよりも低く，$H \ll
\Lambda$ でなければならない．$c_s
\ll 1$ の場合にカットオフの具体的な式 (\ref{eq:cutoff}) を用いると，
\be
H^4 \ll M_{\rm Pl}^2 |\dot H| c_s^5
\ee
を得る．これは，スペクトルの規格化 (\ref{eq:spectrumzetacs}) を用いて，音速に対する不等式
\be
c_s \gg P_\zeta^{1/4} \simeq 0.003 \;. 
\ee
として書き直せる．これよりも音速が小さい理論は，$E \simeq H$ で強結合となる．当然ながら，この音速の値は，非Gaussian性が1程度となる値にも対応する．すなわち，宇宙論的予言に関係するエネルギースケール $H$ において，理論は強結合となるのである．

理論の自然性についてコメントしよう．$c_s$ が小さい極限は自然なのか，それともループ補正によって，より大きな値が誘起されるのかという疑問が生じるかもしれない．例えば $\dot\pi(\nabla\pi)^2$ や $(\nabla\pi)^4$ のようなGoldstoneモードの自己相互作用は，$(\nabla\pi)^2$ に輻射補正を誘起する．これらの寄与が $c_s^{-5}
\Lambda^4/(16 \pi^2 M_2^4)$ 程度であることは容易に見積もれる．ここで $\Lambda$ は紫外カットオフ，{\em すなわち} 新しい物理が加わるエネルギースケールである．大きな輻射補正が生じ得ないことがわかる．$\Lambda$ をユニタリティ限界 (\ref{eq:cutoff}) に取ったとしても，その効果はツリーレベルの値と同程度である．ユニタリティによるカットオフは，まさにループ補正が1程度になるエネルギースケールなので，これは理にかなっている．

また，作用 (\ref{Spi}) は有効場理論の観点から自然であることも指摘しておきたい \cite{Polchinski:1992ed}．実際，$\pi$ の近似的なシフト対称性を仮定すれば，関連演算子は大きな繰り込みから保護される．この場合，関連演算子の係数は十分に小さく，観測にとって重要になることはない．宇宙論的な相関関数が探るのは，$H$ 程度の固定されたエネルギースケールでの理論であり，それよりも低いエネルギーへ進むことはないからである．もちろん，ここではインフレーションの期間だけを考えており，この期間には，近似的なシフト対称性があれば理論を技術的に自然にするのに十分である．インフレーションからの円滑な終了と効率的な再加熱を実現することは，機能する模型に対する追加の要件であり，本形式では論じない．
